\documentclass[a4paper,7pt,onecolumn,oneside]{article}
\usepackage{ctex}
% 或
% \usepackage{xeCJK}
% 这部分声明需要用到的包
\usepackage{amsmath,amsfonts,amssymb,graphicx,bm}    % EPS 图片支持
\usepackage{mathrsfs}
\usepackage{subfigure}   % 使用子图形
\usepackage{amsthm}
\usepackage{enumerate}

\usepackage{geometry}
\usepackage{balance}

\usepackage{indentfirst}

\usepackage{hyperref}


\usepackage{listings}
\usepackage{xcolor}
\lstdefinestyle{ccstyle}{
  language=C++, % 代码语言
  basicstyle=\ttfamily\small, % 代码基本样式
  keywordstyle=\color{blue}, % 关键词颜色
  commentstyle=\color{green!60!black}, % 注释颜色
  stringstyle=\color{red}, % 字符串颜色
  numbers=none, % 行号位置
  numberstyle=\tiny\color{gray}, % 行号样式
  frame=single, % 添加框
  framexleftmargin=5pt, % 框左侧距离代码内容的距离
  backgroundcolor=\color{lightgray!10}, % 代码框背景颜色
  breaklines=true, % 自动换行
  breakatwhitespace=true, % 只在空格断行
  showstringspaces=false, % 不显示字符串中的空格
  captionpos=b % 标题位置：bottom
}

\newtheorem{example}{例}

\begin{document}
LAPACK可用求解器见：
\href{http://sites.science.oregonstate.edu/~landaur/nacphy/lapack/simple.html}{Available DRIVER routines}

函数名中第一个字母代表计算精度：
\begin{itemize}
\item
  s: 代表单精度浮点数（float）。
\item
  d: 代表双精度浮点数（double）。
\item
  c: 代表复数单精度浮点数（complex float）。
\item
  z: 代表复数双精度浮点数（complex double）。
\end{itemize}

\section{变量解释}
\label{sec:variable}
首先进行变量解释，可以结合后面的例子理解
\subsection{\texttt{nrhs}  :  number of right hand side}
\label{sec:nrhs}
\texttt{nrhs} 是一个整数参数，用于指定右手边（Right Hand Side，简称 RHS）的列数。在解线性方程组时，\texttt{nrhs} 指定了右手边的列数，即待求解方程组的个数。

例如，对于一个 n 阶方阵 A，我们想要求解线性方程组 Ax = B，其中 B 是一个 n 行 k 列的矩阵（每列是一个RHS）。那么，在调用 \texttt{LAPACKE\_dgesv} 时，\texttt{nrhs} 应该设置为 k，此时
\[
  B =
  \begin{bmatrix}
    \vert & \vert &  & \vert \\
    b_1 & b_2 & \cdots & b_k \\
    \vert & \vert &  & \vert \\
  \end{bmatrix}.
\]
一般情况下，即 $Ax = b$时，\texttt{nrhs} 为 1，表示只有一个RHS需要求解。
\subsection{\texttt{lda} / \texttt{ldb}  :  leading dimension}
\label{sec:lda}
对于一个 m 行 n 列的矩阵，它通常在内存中以连续的一维数组的形式存储。leading dimension 指定了在这个一维数组中，相邻两行之间的间隔（步长）。

在行主序存储（Row-Major Order）中，leading dimension 就是矩阵的列数 n。每一行的数据在内存中是连续存放的，相邻两行之间的间隔是 n 个元素。

在列主序存储（Column-Major Order）中，leading dimension 就是矩阵的行数 m。每一列的数据在内存中是连续存放的，相邻两列之间的间隔是 m 个元素。

\subsection{\texttt{ipiv}}
\label{sec:ipiv}
\texttt{ipiv} 是一个整数数组，用于存储部分选主元的信息。在线性代数求解中，选主元是高斯消元法中的一种策略，用于避免在计算过程中出现数值不稳定或者误差过大的情况。

在LAPACK中的一些线性方程组求解函数（如 \texttt{LAPACKE\_dgesv}）中，需要提供一个 \texttt{ipiv} 数组作为参数。\texttt{ipiv} 数组的长度应该与矩阵的维度相同（即与矩阵的行数或列数相同），并用于存储选主元的信息。

具体来说，\texttt{ipiv} 数组中存储的是一个排列，例如，如果 \texttt{ipiv[i] = j}，表示在第 \texttt{i} 步的消元过程中，第 \texttt{i} 行和第 \texttt{j} 行进行了交换。


\section{使用一般矩阵求解器}
配置好LAPACK环境后，使用LAPACK的\texttt{dgesv}函数求解3x3线性方程组。考虑 $Ax = b$，其中：
\[
  A =
  \begin{bmatrix}
    2 & 1 & -1 \\
    3 & -2 & 4 \\
    -1 & 3 & -2
  \end{bmatrix},
  \quad
  b =
  \begin{bmatrix}
    1 \\
    4 \\
    3
  \end{bmatrix}.
\]

我们将使用LAPACK的\texttt{dgesv}函数来解决这个方程组。

\subsection{运行成功代码}

使用LAPACK的\texttt{dgesv}函数的代码如下：
\begin{lstlisting}[style=ccstyle]
  #include <iostream>
  #include <lapacke.h>

  int main() {
    double A[9] = {2.0, 1.0, -1.0, 3.0, -2.0, 4.0, -1.0, 3.0, -2.0};
    double b[3] = {1.0, 4.0, 3.0};
    int n = 3; // 矩阵的维度
    int nrhs = 1; // 右手边的列数，即b的列数
    int lda = n; // A的leading dimension，即A矩阵的列数
    int ldb = 1; // B的leading dimension，即B矩阵的列数

    int ipiv[n]; // 存储部分选主元的信息
    int info; // 存储函数的返回值

    // 使用LAPACK的dgesv函数求解方程组
    info = LAPACKE_dgesv(LAPACK_ROW_MAJOR, n, nrhs, A, lda, ipiv, b, ldb);
    if (info != 0) {
      std::cout << "求解方程组失败" << std::endl;
      return 1;
    }

    // 打印解
    std::cout << b[0] << " " << b[1] << " " << b[2] << std::endl;

    return 0;
  }
\end{lstlisting}
可以看出LAPACK把解覆盖在 \texttt{b}。上述矩阵储存的格式为行主序，若要使用列主序需要进行修改，其主要区别在于\texttt{A}和\texttt{B}储存矩阵的方式，为体现 \texttt{B}的区别，将其设为
\[
  B =
  \begin{pmatrix} 
    1  & 2 \\
    4  & 5 \\
    3  & 6
  \end{pmatrix}
\]
\begin{lstlisting}[style=ccstyle]
  double A[3*3] = {2.0, 3.0, -1.0, 1.0, -2.0, 3.0, -1.0, 4.0, -2.0}; // 列主序存储
  double B[3*2] = {1, 4, 3, 2, 5, 6}; // 列主序存储

  int nrhs = 2;
  int lda = n; // A的leading dimension，即A矩阵的行数
  int ldb = n; // B的leading dimension，即B矩阵的行数

  info = LAPACKE_dgesv(LAPACK_COL_MAJOR, n, nrhs, A, lda, ipiv. B, ldb); // 列主序存储
\end{lstlisting}
如果想使用 \texttt{A[][]} 储存数据，有以下例子：
\begin{lstlisting}[style=ccstyle]
  #include <iostream>
  #include <lapacke.h>

  int main() {
    double A[3][3] = {2.0, 1.0, -1.0, 3.0, -2.0, 4.0, -1.0, 3.0, -2.0};
    double b[3] = {1.0, 4.0, 3.0};
    int n = 3; // 矩阵的维度
    int nrhs = 1; // 右手边的列数，即b的列数
    int lda = n; // A的leading dimension，即A矩阵的列数
    int ldb = 1; // B的leading dimension，即B矩阵的列数

    int ipiv[n]; // 存储部分选主元的信息
    int info; // 存储函数的返回值

    // 使用 *A，若前面用 b[][]则使用 *b
    info = LAPACKE_dgesv(LAPACK_ROW_MAJOR, n, nrhs, *A, lda, ipiv, b, ldb);
    if (info != 0) {
      std::cout << "求解方程组失败" << std::endl;
      return 1;
    }

    // 打印解
    std::cout << b[0] << " " << b[1] << " " << b[2] << std::endl;

    return 0;
  }
\end{lstlisting}

如果想使用 \texttt{std::vector} 储存数据，有以下例子：
\begin{lstlisting}[style=ccstyle]
  #include <iostream>
  #include <vector>
  #include <lapacke.h>

  int main() {
    std::vector<double> A = {2.0, 1.0, -1.0, 3.0, -2.0, 4.0, -1.0, 3.0, -2.0}; // 行主序存储
    std::vector<double> B = {1.0, 2.0, 4, 5.0, 3.0, 6.0}; // 两个RHS 行主序存储
    int n = 3; // 矩阵的维度
    int nrhs = 2; // 两个RHS
    int lda = n; // A的leading dimension，即A矩阵的列数
    int ldb = 2; // B的leading dimension，即B矩阵的列数

    std::vector<int> ipiv(n); // 存储部分选主元的信息
    int info; // 存储函数的返回值

    info = LAPACKE_dgesv(LAPACK_ROW_MAJOR, n, nrhs, &A[0], lda, &ipiv[0], &B[0], ldb); // 行主序存储
    if (info != 0) {
      std::cout << "求解方程组失败" << std::endl;
      return 1;
    }

    // 打印解
    for (int i = 0; i < n*ldb; i+=ldb) {
      std::cout << "x" << i/ldb+1 << " = " << B[i] << ", y" << i/ldb+1 << " = " << B[i + 1] << std::endl;
    }

    return 0;
  }
\end{lstlisting}

\section{使用带状矩阵求解器}
以下矩阵中空白部分为 $0$，考虑 $Ax = b$，其中
\[
  A =
  \begin{pmatrix} 
    a_{11}  & a_{12} & a_{13} &  &  &  \\
    a_{21}  & a_{22} & a_{23} & a_{24} &  &  \\
            & a_{32} & a_{33} & a_{34} & a_{35} &  \\
            &  & a_{43} & a_{44} & a_{45} & a_{46} \\
            &  &  & a_{54} & a_{55} & a_{56} \\
            &  &  &  & a_{65} & a_{66}
  \end{pmatrix}
\]
其压缩储存形式为
\[
  AB(KL + KU + 1 + i - j , j) = A(i, j) \quad \max(1,j-KU) \le i \le \min(N,j+KL).
\]
此时 $KL = 1$ 为下对角线数，$KL = 2$ 为上对角线数。
即
\[
  AB =
  \begin{pmatrix} 
    &  &  &  &  &  \\
    &  & a_{13} & a_{24} & a_{35} & a_{46} \\
    & a_{12} & a_{23} & a_{34} & a_{45} & a_{56} \\
    a_{11}  & a_{22} & a_{33} & a_{44} & a_{55} & a_{66} \\
    a_{21}  & a_{32} & a_{43} & a_{54} & a_{65} & 
  \end{pmatrix}.
\]
\section{运行成功代码}

\begin{lstlisting}[style = ccstyle]
  #include <iostream>
  #include <lapacke.h>

  int main() {
    double AB[5][6] ={
      {0, 0, 0, 0, 0, 0},
      {0, 0, 3, 3, 3, 3},
      {0, 2, 2, 2, 2, 2},
      {7, 7, 7, 7, 7, 7},
      {5, 5, 5, 5, 5, 0}
    };
    double b[6] = {1, 4, 3 , 2 , 1 ,3};
    int n = 6; // 矩阵的维度
    int kl = 1;
    int ku = 2;
    int nrhs = 1; // 右手边的列数，即b的列数
    int ldab = n; // A的leading dimension，即A矩阵的列数
    int ldb = 1; // B的leading dimension，即B矩阵的列数

    int ipiv[n]; // 存储部分选主元的信息
    int info; // 存储函数的返回值

    // 使用 *AB，若前面用 b[][]则使用 *b
    info = LAPACKE_dgbsv(LAPACK_ROW_MAJOR, n, kl, ku, nrhs, *AB, ldab, ipiv, b, ldb );
    if (info != 0) {
      std::cout << "求解方程组失败" << std::endl;
      return 1;
    }

    // 打印解
    for (auto i = 0; i < 6; ++i) {
      std::cout << b[i] << "\n";
    }

    return 0;
  }

\end{lstlisting}


\end{document}
